\section{Datasheet for Datasets}
\label{app:datasheet}

This appendix follows the Datasheets for Datasets template
of~\citet{Gebru2018Datasheets}. Project-Utopia is not a static
text/image corpus: the released artifact is a \emph{deterministic
simulator plus six procedurally generated scenarios plus four
verbatim LLM prompts}. Instances are \texttt{(seed, scenario, tick)}
trajectories generated on demand, not pre-materialized rows. We
answer the Gebru template questions under that framing.

\subsection{Motivation}

\begin{itemize}[leftmargin=*]
\item \textbf{For what purpose was the dataset created?}
To evaluate LLM long-horizon hierarchical decision-making and
multi-resource allocation across a fixed 4-channel schema-validated
tool surface. Existing benchmarks evaluate single agents picking
single tools (AgentBench, GAIA, ToolLLM) or flat MARL on stateless
games (MeltingPot, Crafter, BALROG); the orchestration-of-fixed-tools
regime that dominates production agent systems is under-evaluated.
\item \textbf{Who created the dataset and on behalf of which
entity?}
Anonymous authors as an independent academic submission. No
corporate or institutional sponsorship.
\item \textbf{Who funded the dataset's creation?}
Self-funded. All LLM API costs for the 9-experiment matrix
($\sim$2\,250 runs) were paid out of pocket. No grant identifiers.
\item \textbf{Any other comments?}
The benchmark is forked from a public v0.10.0 colony-simulation
game ($\sim$78\,k LOC player-facing surface), stripped across six
refactor stages (S1--S6) to a headless deterministic harness. The
original game code is MIT-licensed and authored by the same
individual; no third-party code with incompatible licenses was
incorporated.
\end{itemize}

\subsection{Composition}

\begin{itemize}[leftmargin=*]
\item \textbf{What do the instances represent?}
Each instance is a simulator trajectory parameterized by
\texttt{(seed, scenario, tick-budget, adapter)}. A trajectory is
the deterministic sequence of states the substrate produces when
ticked under a specific LLM adapter responding to the four channel
prompts. Trajectories are generated on demand, not stored.
\item \textbf{How many instances are there in total?}
Unbounded by construction: 6 scenarios $\times$ $2^{32}$ seeds
$\times$ arbitrary tick budgets. The paper-run matrix
materializes $\sim$2\,250 trajectories (9 experiments
$\times$ 5--10 seeds $\times$ 5 scenarios) on demand from the
released CLI.
\item \textbf{What does the released artifact contain?}
(1) Four verbatim system prompts in \texttt{project\_utopia/data/prompts/}
defining the LLM contract surface;
(2) six hand-tuned scenario blueprints (\S\ref{app:scenarios});
(3) the deterministic simulator core
($\sim$5\,500 LOC Python);
(4) the scoring pipeline (5 dimension plugins + ScoringEngine +
DimensionNormalizer);
(5) two reference adapters (\texttt{ScriptedOraclePolicy},
\texttt{FlatBaselineAdapter}) bracketing the sandwich-norm;
(6) the determinism manifest hashing every Tier-1 state
slice (Appendix~\ref{app:determinism-report});
(7) the NDJSON paper-run output schema (one row per
\texttt{(seed, scenario, tick-window)}) with pydantic validation;
(8) a Docker image pinned to commit
\texttt{fe55b82}.
\item \textbf{Is any information missing from individual
instances?}
No. Each NDJSON row is fully self-describing: schema version,
seed, scenario, tick window, dimension scores, dimension
inputs, adapter ID, and proxy / LLM snapshot ID (if applicable).
\item \textbf{Are there recommended data splits?}
The 9 experiments (E1--E9) define the evaluation cells; within
each, seeds 0--9 are the canonical splits. Held-out seeds 10--19
are reserved for replication and not used in any reported
finding.
\item \textbf{Are there any errors, sources of noise, or
redundancies?}
The substrate is deterministic: any seed produces a single
trajectory bit-identically (Tier 1; \S\ref{sec:reproducibility}).
The only stochasticity above the substrate is the LLM itself; we
isolate this via the three-tier model. No deduplication is
required.
\item \textbf{Does it contain confidential, sensitive, or
personal data?}
No. The benchmark contains no human-subject data, no scraped
web content, no PII, and no copyrighted material beyond
self-authored code.
\end{itemize}

\subsection{Collection Process}

\begin{itemize}[leftmargin=*]
\item \textbf{How was the data acquired?}
Nothing was scraped or licensed. Scenarios were hand-designed
to span a (resource-scarcity $\times$ threat-profile $\times$
terrain-difficulty) coverage matrix; each was independently
validated for guard-idempotency
(\S\ref{app:scenarios}). The four prompts were iteratively
refined across the S1--S6 academic refactor stages, with each
revision logged in
\texttt{docs/ai-research/refactor-plan.md}. The final
prompts are released verbatim with the simulator.
\item \textbf{Who was involved in data collection?}
The single author. No crowdworkers, no external annotators,
no human subjects.
\item \textbf{Over what timeframe was the data collected?}
The scenario blueprints and prompts were authored between
2025-12 and 2026-04. The deterministic substrate was forked
from the v0.10.0 game (tag \texttt{pre-academic-refactor-v0.10.0},
commit \texttt{16a593d}).
\item \textbf{Were any ethical review processes conducted?}
No IRB review was required: no human subjects, no
behavioral data, no PII.
\item \textbf{Was consent obtained?}
N/A (no data collected from people).
\end{itemize}

\subsection{Preprocessing, Cleaning, and Labeling}

\begin{itemize}[leftmargin=*]
\item \textbf{Was any preprocessing done?}
None for prompts or scenarios (released verbatim from their
authoring repo). On the output side, every NDJSON row is
validated against a pydantic schema and every numeric
directive emitted by an LLM is clamped through
\texttt{project\_utopia.simulation.ai.llm.guardrails}
before being forwarded to the substrate
(weights $\in [0,3]$, risk $\in [0,1]$, TTL $\in$
[8\,h, 24\,h]; full table in Appendix~\ref{app:prompts}).
\item \textbf{Was the raw data saved?}
Yes. The unguarded LLM response (verbatim JSON text and
guard-clamping deltas) is preserved per row in
\texttt{response\_raw} alongside the validated
\texttt{response\_clamped}, enabling ex-post analysis of
schema-violation patterns.
\item \textbf{Is the preprocessing software available?}
Yes: pydantic schemas
(\texttt{project\_utopia.simulation.ai.llm.response\_schema}),
guard-clamping rules
(\texttt{project\_utopia.simulation.ai.llm.guardrails}), and the
5-dimension plugin code are all in the release tree under
Apache 2.0 / MIT.
\end{itemize}

\subsection{Uses}

\begin{itemize}[leftmargin=*]
\item \textbf{Has the dataset been used for any tasks already?}
The 9-experiment paper-run matrix in this submission
(\S\ref{sec:experiments}). No external uses to date.
\item \textbf{What other tasks could this be used for?}
Evaluating (i) hierarchical decision-making across a fixed
tool surface; (ii) cross-vendor channel routing; (iii)
memory degradation under stale long-horizon context;
(iv) reproducibility-aware LLM evaluation methodology
(LayerCast, RecordReplayCache); (v) RLHF / DPO with the
4-channel directive surface as the action space (future
work, \S\ref{sec:limitations}).
\item \textbf{Is there anything that should not be done with
this dataset?}
The prompts and scenarios should \textbf{not} be used as
LLM pre-training or fine-tuning corpora. Doing so would
contaminate any subsequent capability claim. The
verbatim-release-as-action-space design assumes the prompts
remain outside training distributions; we ask downstream
users to honor this contract.
The benchmark is also \emph{not} a substitute for:
single-task tool-call accuracy benchmarks (use AgentBench,
GAIA, ToolBench); long-context language QA (use LongBench,
$\infty$Bench); or general-knowledge / reasoning capability
batteries (use MMLU, GPQA, HumanEval). Project-Utopia tests
hierarchical decision-making under schema-validated
long-horizon resource pressure, nothing broader
(\S\ref{sec:limitations}).
\item \textbf{Any tasks for which the dataset is unsuitable?}
Multimodal perception, embodied manipulation, free-form
linguistic dialogue, and open-world tool discovery are all
out of scope. The action space is fixed at 4 channels by
design; benchmarks of tool-\emph{selection} quality should
look elsewhere.
\end{itemize}

\subsection{Distribution}

\begin{itemize}[leftmargin=*]
\item \textbf{Will the dataset be distributed to third parties?}
Yes, publicly. The source tree is on GitHub at
\texttt{[ANONYMIZED-REPO-URL]} (de-anonymized for camera-ready);
the installable package is on PyPI as \texttt{project-utopia}.
\item \textbf{How will it be distributed?}
\texttt{pip install project-utopia}; the source tree is also
checkout-able at git tag
\texttt{refactor/academic-benchmark-py-rc1}. A Docker image
pinned to commit \texttt{fe55b82} is published alongside.
\item \textbf{When will it be distributed?}
Pre-print release coincides with paper submission
(2026-06-07 NeurIPS D\&B deadline). The repository has been
public throughout the refactor.
\item \textbf{Under what license?}
Apache 2.0 for the simulator core and benchmark framework;
MIT for the prompts and scenario blueprints. All artifacts
are open and freely re-distributable. No proprietary data,
no third-party data, no copyright restrictions.
\item \textbf{Are there fees or access restrictions?}
None.
\end{itemize}

\subsection{Maintenance}

\begin{itemize}[leftmargin=*]
\item \textbf{Who will support, host, and maintain the dataset?}
The author, via the GitHub repository above. The PyPI
package is auto-published from the same source tree on tag.
\item \textbf{How can the owner be contacted?}
Via GitHub issues or the email on the de-anonymized
repository.
\item \textbf{Is there an erratum?}
Yes: \texttt{CHANGELOG.md} in the repository root tracks
every release. Datasheet revisions are stamped under the
current refactor section heading.
\item \textbf{Will the dataset be updated?}
Yes. Versions are pinned via git tags
(\texttt{refactor/academic-benchmark-py-rc1}, etc.). Reviewers
who wish to replicate the exact paper-run matrix should
check out the tagged commit. Subsequent minor releases
(e.g., scenario additions) will be tagged separately and
will not retroactively change paper-reported numbers.
\item \textbf{Will older versions remain available?}
Yes. Every tagged version remains on GitHub and PyPI
indefinitely.
\item \textbf{Can others extend / build on the dataset?}
Yes. The \texttt{AgentAdapter} contract is the documented
extension surface for new LLM backends; new scenarios
plug into \texttt{ScenarioSampler}; new dimensions
implement the \texttt{DimensionPlugin} protocol. We
encourage forks and pull requests but do not require
attribution beyond the license.
\end{itemize}
